% =====================================================================
\chapter{Distribuciones muestrales y teorema central del límite}
\label{tema:3}
% =====================================================================

Un estadístico es una variable aleatoria y, para hacer inferencia, hace falta conocer (al menos aproximadamente) su distribución. Este tema estudia la distribución de la media muestral: exacta cuando los datos son normales y aproximada, gracias al teorema central del límite, en cualquier otro caso. Se incluyen la aproximación normal a la binomial y las distribuciones $\chi^{2}$ y $t$ de Student, que se usarán en los intervalos de confianza y en los contrastes. Se sigue el capítulo~10 de Blitzstein y Hwang~\cite{bh} (§10.3 «Central limit theorem» y §10.4 «Chi-Square and Student-t»), la lectura~6b del curso 18.05 del MIT~\cite{mit1805}, «Central Limit Theorem and the Law of Large Numbers», y el capítulo~7 de OpenStax~\cite{openstax} (§7.1 y §7.2).

\section{La distribución normal}

\begin{definicion}[Distribución normal]\label{def:t3-normal}
Una variable $X$ tiene distribución \emph{normal} de media $\mu\in\R$ y varianza $\sigma^{2}>0$, y se escribe $X\sim\Nor(\mu,\sigma^{2})$, si su densidad es
\[
f(x)=\frac{1}{\sigma\sqrt{2\pi}}\exp\Bigl(-\frac{(x-\mu)^{2}}{2\sigma^{2}}\Bigr),\qquad x\in\R.
\]
La \emph{normal estándar} es $Z\sim\Nor(0,1)$; su función de distribución se denota $\Phi(z)=\PP(Z\le z)$ y su cuantil de orden $q$, $z_q=\Phi^{-1}(q)$.
\end{definicion}

La función $\Phi$ no tiene expresión mediante funciones elementales y se tabula (\cref{tab:phi}). Por la simetría de la densidad respecto de~$0$, $\Phi(-z)=1-\Phi(z)$ y $z_{1-q}=-z_q$.

\begin{proposicion}[Tipificación]\label{prop:t3-tipificacion}
Si $X\sim\Nor(\mu,\sigma^{2})$ y $a\neq0$, entonces $aX+b\sim\Nor(a\mu+b,a^{2}\sigma^{2})$. En particular, $Z=(X-\mu)/\sigma\sim\Nor(0,1)$ y
\[
\PP(X\le x)=\Phi\Bigl(\frac{x-\mu}{\sigma}\Bigr).
\]
\end{proposicion}

La demostración es el \cref{eje:t3-lineal-normal}. Con la tabla se obtiene la conocida regla empírica: una normal toma valores a menos de $1$, $2$ y $3$ desviaciones típicas de su media con probabilidades $2\Phi(1)-1\approx0{,}683$, $2\Phi(2)-1\approx0{,}954$ y $2\Phi(3)-1\approx0{,}997$.

\begin{teorema}[Combinaciones lineales de normales independientes]\label{teo:t3-suma-normales}
Si $X_1,\dots,X_n$ son independientes con $X_i\sim\Nor(\mu_i,\sigma_i^{2})$, entonces $\sum_ia_iX_i\sim\Nor\bigl(\sum_ia_i\mu_i,\sum_ia_i^{2}\sigma_i^{2}\bigr)$ para cualesquiera $a_i$ no todos nulos. En particular, si $X_1,\dots,X_n\sim\Nor(\mu,\sigma^{2})$ son \iid,
\[
\Xbar\sim\Nor\Bigl(\mu,\frac{\sigma^{2}}{n}\Bigr).
\]
\end{teorema}

La media y la varianza se siguen de la \cref{prop:t2-props}; que la distribución siga siendo normal se demuestra con funciones generadoras de momentos~\cite[§6.6]{bh} y se admite aquí sin demostración.

\section{El teorema central del límite}

Cuando los datos no son normales, la distribución exacta de $\Xbar$ suele ser complicada. El teorema central del límite afirma que, sin embargo, para $n$ grande es aproximadamente normal, sea cual sea la distribución de partida, con tal de que tenga varianza finita.

\begin{teorema}[Teorema central del límite]\label{teo:t3-tcl}
Sean $X_1,X_2,\dots$ \iid\ con media $\mu$ y varianza $0<\sigma^{2}<\infty$, y sea $S_n=X_1+\dots+X_n$. Entonces, para todo $z\in\R$,
\[
\lim_{n\to\infty}\PP\Bigl(\frac{\Xbar_n-\mu}{\sigma/\sqrt n}\le z\Bigr)=\lim_{n\to\infty}\PP\Bigl(\frac{S_n-n\mu}{\sigma\sqrt n}\le z\Bigr)=\Phi(z).
\]
\end{teorema}

Las dos variables tipificadas son la misma (basta multiplicar numerador y denominador por $n$). La demostración, con funciones generadoras de momentos, está en~\cite[§10.3]{bh}; aquí interesa su uso.

\begin{observacion}[Uso práctico]\label{obs:t3-uso}
Para $n$ «grande» se aproxima
\[
\Xbar\approx\Nor\Bigl(\mu,\frac{\sigma^{2}}{n}\Bigr),\qquad S_n\approx\Nor(n\mu,n\sigma^{2}).
\]
Una regla habitual es $n\ge30$~\cite[§7.1]{openstax}, pero la calidad de la aproximación depende de la forma de la distribución: con datos simétricos basta un $n$ pequeño y con datos muy asimétricos hace falta bastante más. El teorema exige varianza finita: para distribuciones de cola muy pesada, como la de Cauchy (o la $t$ con un grado de libertad), la media muestral no se estabiliza~\cite[§10.3]{bh}.
\end{observacion}

\addtocounter{definicion}{1}%

\section{Aproximación normal a la binomial}

Si $X\sim\Bin(n,p)$, entonces $X$ tiene la misma distribución que la suma de $n$ variables $\Ber(p)$ \iid, con $\mu=p$ y $\sigma^{2}=p(1-p)$. El teorema central del límite da el resultado clásico de De Moivre y Laplace.

\begin{corolario}[Aproximación normal a la binomial]\label{cor:t3-binomial}
Si $X\sim\Bin(n,p)$ con $0<p<1$, para $n$ grande
\[
X\approx\Nor\bigl(np,\,np(1-p)\bigr),\qquad \widehat p=\frac Xn\approx\Nor\Bigl(p,\,\frac{p(1-p)}{n}\Bigr).
\]
Una regla habitual para usarla es $np\ge10$ y $n(1-p)\ge10$.
\end{corolario}

\begin{definicion}[Corrección por continuidad]\label{def:t3-cc}
Al aproximar una variable con valores enteros por una continua, cada valor entero $k$ se identifica con el intervalo $[k-\frac12,k+\frac12]$. Así, para enteros $a\le b$,
\[
\PP(a\le X\le b)\approx\Phi\Bigl(\frac{b+\frac12-np}{\sqrt{np(1-p)}}\Bigr)-\Phi\Bigl(\frac{a-\frac12-np}{\sqrt{np(1-p)}}\Bigr).
\]
\end{definicion}

\addtocounter{definicion}{1}%

\section{Distribuciones \texorpdfstring{$\chi^{2}$}{chi cuadrado} y \texorpdfstring{$t$}{t} de Student}

Cuando $\sigma$ es desconocida y se sustituye por $S$, la variable $(\Xbar-\mu)/(S/\sqrt n)$ ya no es normal. Su distribución exacta, para datos normales, se describe con dos familias derivadas de la normal~\cite[§10.4]{bh}.

\begin{definicion}[Distribuciones $\chi^{2}$ y $t$]\label{def:t3-chi-t}
Sean $Z_1,\dots,Z_k\sim\Nor(0,1)$ \iid\ La distribución de $V=Z_1^{2}+\dots+Z_k^{2}$ se llama \emph{$\chi^{2}$ con $k$ grados de libertad} y se escribe $V\sim\chi^{2}_k$. Si $Z\sim\Nor(0,1)$ y $V\sim\chi^{2}_k$ son independientes, la distribución de
\[
T=\frac{Z}{\sqrt{V/k}}
\]
se llama \emph{$t$ de Student con $k$ grados de libertad} y se escribe $T\sim t_k$. Sus cuantiles se denotan $\chi^{2}_{k;q}$ y $t_{k;q}$ (\cref{tab:t,tab:chi2}).
\end{definicion}

La $\chi^{2}_k$ toma valores positivos, tiene media $k$ y varianza $2k$ (\cref{eje:t3-chi-momentos}). La $t_k$ es simétrica respecto de~$0$ (\cref{eje:t3-t-simetrica}), tiene colas más pesadas que la normal y converge a ella cuando $k\to\infty$: por ejemplo, $t_{5;0{,}975}=2{,}571$, $t_{30;0{,}975}=2{,}042$ y $z_{0{,}975}=1{,}960$.

\begin{teorema}[Teorema de Fisher]\label{teo:t3-fisher}
Sean $X_1,\dots,X_n\sim\Nor(\mu,\sigma^{2})$ \iid\ con $n\ge2$. Entonces $\Xbar$ y $S^{2}$ son independientes, $\Xbar\sim\Nor(\mu,\sigma^{2}/n)$ y
\[
\frac{(n-1)S^{2}}{\sigma^{2}}\sim\chi^{2}_{n-1}.
\]
\end{teorema}

Se admite sin demostración~\cite[§10.4]{bh}. Los $n-1$ grados de libertad son los de la \cref{obs:t1-n-1}: la suma de cuadrados respecto de $\Xbar$ equivale a una suma de $n-1$ cuadrados de normales estándar independientes. Como $\E[\chi^{2}_{n-1}]=n-1$, el teorema reproduce, para datos normales, el resultado $\E[S^{2}]=\sigma^{2}$ del \cref{teo:t2-s2}, y como $\Var(\chi^{2}_{n-1})=2(n-1)$, da $\Var(S^{2})=2\sigma^{4}/(n-1)$, que se usó en el \cref{eje:t2-ecm-normal}.

\begin{corolario}[Estadístico $t$]\label{cor:t3-t}
En las condiciones del \cref{teo:t3-fisher},
\[
T=\frac{\Xbar-\mu}{S/\sqrt n}\sim t_{n-1}.
\]
\end{corolario}

\addtocounter{definicion}{1}%

\addtocounter{definicion}{1}%

